\documentclass[../../main.tex]{subfiles}
\begin{document}

\begin{center}
  {\fontsize{48pt}{58pt}\selectfont \textbf{1996}\normalsize\textbf{年度}}
\end{center}

\begin{table}[htb]
  \centering
  \caption{出題テーマと難易度}
  \begin{tabular}{|c|p{9cm}|c|}
    \hline
    \textbf{問題番号} & \textbf{テーマ} & \textbf{難易度} \\
    \hline
    第1問 & 三角関数を表す整数係数多項式 & \\
    \hline
    第2問 & 累乗数の末尾と10の倍数 & \\
    \hline
    第3問 & 2直線に関する鏡映と点列 & \\
    \hline
    第4問 & 四面体の辺の長さと点の存在範囲 & \\
    \hline
    第5問 & 正三角形の移動でできる立体 & \\
    \hline
    第6問 & 円周上の確率的な移動 & \\
    \hline
  \end{tabular}
\end{table}

\begin{table}[htb]
  \centering
  \caption{解答}
  \renewcommand{\arraystretch}{1.4}
  \begin{tabular}{|c|c|p{9cm}|}
    \hline
    \textbf{問題番号} & \textbf{小問} & \textbf{答え} \\
    \hline
    第1問 & (1)〜(3) & 証明問題 \\
    \hline
    第2問 & -- & $(m,n)=(9,19)$の一例 \\
    \hline
    第3問 & (1) & 証明問題 \\
    \cline{2-3}
    & (2) & 最大値$8$ \\
    \hline
    第4問 & -- & $x+y>0$，$x-y>0$，$\dfrac{(x-2)^2}{3}+y^2<1$ \\
    \hline
    第5問 & (1) & $V=\dfrac{\sqrt{3}\pi}{3}(3a-2a^3)$ \\
    \cline{2-3}
    & (2) & $a=\dfrac{1}{\sqrt{2}}$で最大値$\dfrac{\sqrt{6}\pi}{3}$ \\
    \hline
    第6問 & (1) & $p_n=\dfrac{1}{3}\left(1-\left(-\dfrac{1}{2}\right)^n\right)$ \\
    \cline{2-3}
    & (2) & $2\left(\dfrac{1}{3}\right)^k$ \\
    \hline
  \end{tabular}
\end{table}

\paragraph{試験形態}
6問．試験時間は要確認．

\end{document}
